% Reduced extract from Bos and McCurley, arXiv:2301.08277v2.
% License: CC BY 4.0. See SOURCES.md for attribution and source archive.
\documentclass[version=preprint]{iacrcc}
\newcommand{\cmd}[1]{\texttt{\textbackslash{}#1}}
\title[subtitle={(updated April 3, 2023)}]{LaTeX, metadata, and publishing workflows}
\addauthor[orcid={0000-0003-1010-8157},inst={1}]{Joppe W. Bos}
\addauthor[orcid={0000-0001-7890-5430},inst={2}]{Kevin S. McCurley}
\affiliation[country={Belgium}]{NXP Semiconductors}
\affiliation[country={United States}]{Self}

\begin{document}
\maketitle
\keywords[publishing, LaTeX, metadata]{publishing, \LaTeX, metadata}

\begin{abstract}
The field of scientific publishing that is served by \LaTeX\ is increasingly
dependent on the availability of metadata about publications. We discuss how
to use \LaTeX\ classes and \BibTeX\ styles to curate metadata throughout the
life cycle of a published article. Our focus is on streamlining and automating
much of publishing workflow.
\end{abstract}

\section{Introduction}
The original goal of \TeX\ was focused primarily on typesetting, and the
appearance of the output on paper. The later invention of \LaTeX\ was focused
on "letting the user concentrate on the structure of the text rather than on
formatting commands". Users were encouraged to write their papers using
high-level macros like \texttt{\textbackslash{}section}, and leave decisions
about appearance to the style that is used.

There is at least one area in which the \LaTeX\ community has been slow to
adapt to the needs of modern publishing workflows, namely in the
{\em curation of metadata about publications}. This is the main focus of this
article. Some of this metadata is fairly obvious, such as the title, list of
authors, affiliations, funding sources, and references.
\end{document}
